\section{Resultados}\label{sec:dar-resultados}

A partir de la aplicación del análisis \ac{DAR} a las alternativas identificadas, se obtuvo una valoración comparativa basada en los criterios de selección definidos para el estudio. Estos criterios permitieron determinar el nivel de correspondencia de cada alternativa con las características consideradas relevantes para la selección de una tecnología de servicio de directorio.

Los resultados mostraron diferencias entre las alternativas evaluadas en aspectos relacionados con sus capacidades de autenticación y gestión de identidades, mecanismos de seguridad, interoperabilidad, administración, documentación y soporte, entre otros criterios establecidos. La valoración conjunta de estos aspectos permitió establecer una alternativa con un mayor nivel de correspondencia frente al conjunto de criterios definidos.

Como resultado del análisis, FreeIPA fue la tecnología que obtuvo la valoración más favorable, por lo que fue seleccionada como la alternativa para continuar con las siguientes fases del proyecto. Esta selección fue resultado de la evaluación sistemática de las alternativas mediante los criterios establecidos y no de la consideración aislada de una característica particular de las tecnologías analizadas.

De esta manera, el análisis \ac{DAR} permitió resolver la selección de la tecnología a utilizar, estableciendo a FreeIPA como la alternativa seleccionada para la implementación del servicio de directorio.
